\booksection{第一章附：鴻蒙操作系統專題}

\chapterintro{鴻蒙（HarmonyOS）乃華爲自主研發之操作系統，自2019年發佈至今，已歷數個大版本。其不唯“又一安卓替代品”，乃有獨特設計理念之新物種。讓吾等一同深入識此正在崛起之操作系統。}

\booksubsection{一、鴻蒙爲何物？何以與衆不同？}

鴻蒙之核心理念乃“萬物互聯”。異於傳統操作系統主要服務於單一設備，鴻蒙自設計之初即考慮了多設備協同。

\textbf{分佈式架構}乃鴻蒙之靈魂。簡言之，其可將多個設備虛擬成一個“超級終端”。汝之手機、平板、計算機、手錶、電視，不再是孤立之個體，乃一個整體之“器官”，可互相調用對方之能力。

譬如：汝用手機觀視頻，欲於大屏上觀，傳統方式乃投屏；而於鴻蒙上，汝可直接將視頻“流轉”至電視上續觀，手機化爲遙控器。此非簡單之投屏，乃應用本身於跨設備運行。

\booksubsection{二、鴻蒙4.3：基于AOSP之成熟版本}

鴻蒙4.3乃當前華爲手機上廣泛部署之版本。須說明者，鴻蒙4.3兼容安卓開源項目（AOSP）之代碼框架，故其能完整兼容安卓應用——汝於華爲應用市場下載之大部分App，本質上仍是安卓應用，通過鴻蒙之兼容層運行。然鴻蒙4.3並非簡單套殼安卓，其於AOSP基礎上加入了大量華爲自研之分佈式能力、方舟引擎等技術。

\subsubsection*{核心特性：}

\begin{itemize}
  \item \textbf{安卓應用兼容}：由於基於AOSP開發，鴻蒙4.3可正常運行幾乎所有安卓應用，用戶幾乎感受不到差異。
  \item \textbf{方舟引擎}：華爲自研之編譯器與運行時，使應用運行更流暢、更省電。較早期版本，流暢度提升明顯。
  \item \textbf{超級終端}：於控制中心，可直觀地見周圍之其他鴻蒙設備，拖拽即可連接。手機調用計算機之鍵盤鼠標、平板調用手機之攝像頭，皆甚方便。
  \item \textbf{服務卡片}：類似iOS之小組件，然更靈活。可上滑App圖標展開卡片，顯示關鍵信息，不用打開App。
  \item \textbf{純淨之模}：唯允許從官方應用市場安裝應用，減惡意軟件風險。
  \item \textbf{隱私中心}：集中管理應用權限，觀應用行爲記錄。
\end{itemize}

\subsubsection*{用戶界面特點：}

\begin{itemize}
  \item \textbf{雜誌鎖屏}：鎖屏界面顯示精美圖片與文章，拂去即入桌面。
  \item \textbf{萬能卡片}：桌面圖標可展開成卡片，顯示更多信息。
  \item \textbf{超級中轉站}：一個臨時存儲區，可將文字、圖片、文件暫存，而後拖至任意App中用。跨應用分享變得無比絲滑。
  \item \textbf{智慧搜索}：全局檢索，可搜應用、文件、設置、網頁。
\end{itemize}

\booksubsection{三、鴻蒙NEXT（鴻蒙6）：華爲自主操作系統之真正起點}

鴻蒙NEXT（HarmonyOS NEXT 5.0）與鴻蒙4.3有根本性之別——其採華爲自研之鴻蒙微內核，不再依賴Linux內核與AOSP代碼，亦\textbf{不再直接支持安卓應用}。此乃鴻蒙自“兼容安卓之過渡形態”正式走向“完全獨立之操作系統”之關鍵一步。

\subsubsection*{重大變革：}

\begin{itemize}
  \item \textbf{完全自主內核}：鴻蒙NEXT用華爲自研之內核，不再包含任何AOSP代碼。系統更精簡、更安全、更高效。
  \item \textbf{不再支持安卓應用}：此乃最大之變化。鴻蒙NEXT唯能運行原生鴻蒙應用，無法安裝與運行APK文件。升級前請確認汝常用之App已有鴻蒙原生版本。
  \item \textbf{原生應用生態}：華爲正全力推動開發者將應用遷移至鴻蒙原生平臺。目前主流App如微信、支付寶、淘寶等已推出鴻蒙原生版本，然部分小衆應用或尚未適配。
  \item \textbf{星盾安全架構}：全新之安全體系，自硬件至軟件之全鏈路保護。支持更細粒度之權限控制，譬如"模糊定位"——應用唯知汝大概何處，不知精確位置。
  \item \textbf{AI深度融合}：系統級AI能力，不唯某個App有AI。譬如AI字幕可實時翻譯視頻，AI摳圖可智能識別主體。
  \item \textbf{跨設備無縫流轉}：更進一步之分佈式能力。應用狀態可於設備間無縫切換，汝於手機上編輯之文檔，行至計算機前自動於計算機上打開。
  \item \textbf{元服務}：無需安裝即可用之輕量級服務。譬如掃二維碼直接打開某個功能，不用先下載App。
\end{itemize}

\begin{tcbwarn}{升級鴻蒙NEXT前必讀}
升級鴻蒙NEXT前必讀

\end{tcbwarn}

\subsubsection*{操作邏輯變化：}

\begin{itemize}
  \item \textbf{手勢導航}：左右邊緣滑動返回，底部上滑回桌面，上滑停頓進多任務。
  \item \textbf{智慧多窗}：支持分屏、懸浮窗、平行視界等多種多任務形態。
  \item \textbf{暢連通話}：華爲之視頻通話服務，支持屏幕共享、美顏、背景虛化。
\end{itemize}

\booksubsection{四、鴻蒙生態：設備與應用}

鴻蒙不唯手機系統，乃一個覆蓋多種設備之生態系統。

\subsubsection*{支持之設備類型：}

\begin{itemize}
  \item \textbf{手機與平板}：核心設備，功能最完整。
  \item \textbf{智能手錶與手環}：獨立運行或與手機協同。
  \item \textbf{智慧屏（電視）}：家庭娛樂中心。
  \item \textbf{智能音箱}：語音交互入口。
  \item \textbf{汽車}：鴻蒙座艙，智能駕駛艙體驗。
  \item \textbf{智能之家}：冰箱、空調、洗衣機等。
\end{itemize}

\subsubsection*{應用生態現狀：}

鴻蒙應用市場已有數百萬應用，主流App基本皆有。華爲亦於積極推動開發者開發原生鴻蒙應用。截至2026年中，微信、支付寶、淘寶、抖音、美團、京東等頭部應用均已推出鴻蒙原生版本，日常使用基本無障礙。部分小衆應用與行業應用或仍於適配中，升級前建議確認汝常用之App是否已支持。

\booksubsection{五、鴻蒙實用技巧}

\begin{itemize}
  \item \textbf{超級終端}：自屏幕右上角下滑啓控制中心，點"超級終端"，可見周圍之其他設備，拖拽連接。
  \item \textbf{超級中轉站}：長按文字/圖片/文件，擇"複製到中轉站"，而後於其他App中長按粘貼，或拖拽出來。
  \item \textbf{服務卡片}：上滑App圖標展開卡片，長按卡片可調大小或添至桌面。
  \item \textbf{智慧多窗}：自屏幕左右邊緣向內滑動並停頓，啓側邊欄，點應用以懸浮窗啓。
  \item \textbf{隱私保護}：設置 > 隱私，可觀應用權限使用記錄，設模糊定位等。
\end{itemize}

\begin{tcbwarn}{注意之事}
注意事項

\end{tcbwarn}

\begin{tcbpractice}{今即試之：體驗鴻蒙特色功能}
\begin{enumerate}
  \item 開控制中心，尋"超級終端"，觀四周有無他鴻蒙之設備。
  \item 開控制中心，尋"超級終端"，觀四周有無他鴻蒙之設備。
  \item 上滑數個常用App圖標，觀其服務卡片所顯何內容。
  \item 往設置 > 隱私，查閱近時何應用曾用敏感權限。
\end{enumerate}

\end{tcbpractice}
